% Clase 20 — Dorantes Mejía Emmanuel (potenciómetro / ADC)
\paragraph{Potenciómetro (ADC)}
\textbf{Autor:} Dorantes Mejía Emmanuel

Cada muestra del ADC del RP2040 es de 12 bits (0 a 4095) y se guarda en \texttt{uint16\_t axis[4]}, el entero más
chico que la contiene. El DMA copia con \texttt{DMA\_SIZE\_16} (\texttt{read\_adc.c}, línea 19), así que cada muestra
cae completa en un solo elemento y nunca queda a medias. Se muestrean 4 canales en \textit{round-robin}
(\texttt{read\_adc.c}, línea 46) aunque solo hay 3 ejes: el cuarto (ADC3) completa un anillo de
$4 \times 2 = 8 = 2^3$ bytes, porque el anillo del DMA (\texttt{read\_adc.c}, línea 22) solo acepta tamaños
potencia de 2. Por eso \texttt{g\_controls} lleva \texttt{\_\_attribute\_\_((aligned(8)))} (\texttt{main.c}, línea 24):
el anillo exige que su dirección sea múltiplo de 8.

La frecuencia es de 4000 muestras por segundo por canal, o sea 16\,000 conversiones por segundo en total. Como el
ADC trabaja con un reloj de 48 MHz, el divisor es $48\,000\,000 / 16\,000 - 1 = 2999$. El potenciómetro se actualiza
cada $1/4000 = 0.25$ ms, mucho más rápido que los 10 ms comprometidos en el Avance 1.

El ADC no tiene variable \texttt{volatile} propia: escribe directo en \texttt{g\_controls.axis}, que ya es
\texttt{volatile} por la misma razón que los botones (la escribe el DMA, no la CPU).
